%! TeX root: ../main.tex
\chapter{Nonlinear Conservation Laws in the Geometric Framework}
\label{sec:nonlinear-conservation-laws}

The previous chapters established a robust geometric framework for linear advection.
We derived high-resolution implicit schemes via the Implicit Upwind Range Condition (IURC).
However, many physical phenomena are governed by nonlinear conservation laws where the speed of information propagation depends on the solution itself.
In this chapter, we extend the implicit geometric framework to scalar nonlinear conservation laws of the form \(u_t + f(u)_x = 0\).
The transition from linear to nonlinear equations introduces new phenomena: characteristics are no longer parallel, shocks and expansion waves can form.
The numerical solution can converge to more than one solution: our scheme should choose the physically "right" one, e.g., one that satisfy a numerical entropy condition.

We begin by examining the nonlinear analogue of the first-order implicit upwind scheme, particularly the formulation proposed by Lozano and Aslam \cite{lozano2020implicit}.
Building on this, we derive second-order extensions using geometric arguments.
To stabilize the resulting schemes, we extend IURC to ensure monotonicity and stability in the presence of shocks and rarefactions, using Burgers' equation as our primary laboratory.
Finally, we discuss the numerical implementation via sweeping algorithms.

\section{Nonlinear Fluxes}
\label{sec:nonlinear_fluxes}

The linear advection equation
\begin{equation}
    u_t + a u_x = 0
\end{equation}
assumes a constant wave speed \(a\).
For a general nonlinear conservation law,
\begin{equation}
    \label{eq:nonlinear_cons}
    u_t + f(u)_x = 0,
\end{equation}
rewritten in quasi-linear form as
\begin{equation}
    \label{eq:claw_quasi_lin}
    u_t + f'(u)\,u_x = 0,
\end{equation}
the local wave speed is given by the derivative of the flux function, \(a(u) = f'(u)\).
This nonlinearity implies that different parts of the solution travel at different speeds, leading to the formation of shocks (where characteristics converge) and rarefactions (where characteristics diverge).

To explore our ideas, we will primarily focus on the inviscid Burgers' equation,
\begin{equation}
    \label{eq:burgers}
    u_t + \left( \frac{u^2}{2} \right)_x = 0,
\end{equation}
with flux \( f(u) = \frac{u^2}{2} \) and wave speed \(f'(u) = a(u) = u\).
This equation serves as the simplest model that captures the essential features of nonlinear transport and shock formation.

\section{Geometric Formulation of Implicit Nonlinear Fluxes}%
\label{sec:geometric-implicit-nonlinear}

In this section, we extend our implicit geometric framework to nonlinear implicit schemes, following the logic of Goodman and Leveque \cite{goodman1988geometric}.

% TODO: fig of linear reconstruction u(x), don't forget u_i^+
% TODO: fig of flux approximation f(u)

Our goal is to approximate the time-averaged flux at the interface: 
\begin{equation}
    F_{i+1/2} = \frac{1}{\Delta t} \int\limits_{t^{n}}^{t^{n+1}} f(u(x_{i+1/2},t)) \,\mathrm{d}t.
\end{equation}

Assume \(f(u)\) is a nonlinear, convex flux such that \(f''(u) \geq 0\).
For the following derivation, we assume that the wave speed, thus in our case both \(f'(u_{i}^{n+1})\) and \(a_i\) are positive, where the information propagates from left to right.
To make the integration tractable, we approximate the flux function in the vicinity of the \(i\)-th cell with a straight line:
\begin{equation}
    \label{eq:flux_approx}
    f(u) \approx g_i(u) = f(u_{i}^{n+1}) + a_i \left( u - u_{i}^{n+1} \right),
\end{equation}
where \(a_i\) is an average slope of the flux \(g_i' = a_i\).

Next, we consider a piecewise linear reconstruction of the solution at time \(t^n\) as before:
\begin{equation}
    \tilde{u}(x, t^n) = u_i^{n+1} + \sigma_i(x - (x_i - a\Delta t)),
\end{equation}
where \(\sigma_i\) is the slope in cell \(i\).
For the interface \(x_{i+1/2}\) at any time \(t \in [t^{n}, t^{n+1}]\), the value is obtained by tracing along a characteristic.
Assuming the wave travels locally at a constant speed \(a_i\), the value is given by:
\begin{equation}
    \label{eq:interface_val_approx}
    u(x_{i+1/2}, t) = u_i^+ - \sigma_i a_i \left( t - t^n \right),
\end{equation}
where \(u_i^+ = u_{i}^{n+1} + \sigma_i \left( f'(u_{i}^{n+1})\Delta t + \frac{\Delta x}{2} \right)\) represents the value at the right interface \(x_{i+1/2}\) at time \(t^n\).
Notice that here we assume that the characteristic speed in cell \(i\) at time \(t^{n+1}\) is \(f'(u_{i}^{n+1})\), more on that later.

By substituting (\ref{eq:interface_val_approx}) into the linearized flux (\ref{eq:flux_approx}), we obtain the integrand for the numerical flux:
\begin{equation}
    \begin{split}
        f(u(x_{i+1/2},t)) \approx g_i \left( u(x_{i+1/2}, t) \right)
        &= f(u_{i}^{n+1}) + a_i \left( u(x_{i+1/2}, t) - u_{i}^{n+1} \right) \\
        &= f(u_{i}^{n+1}) + a_i \left( u_i^+ - \sigma_i a_i \left( t - t^n \right) - u_{i}^{n+1} \right) \\
        &= f(u_{i}^{n+1}) + a_i \left( u_i^+ - u_{i}^{n+1} - \sigma_i a_i \left( t - t^n \right) \right).
    \end{split}
\end{equation}
Integrating the above expression over the time step \([t^{n}, t^{n+1}]\) yields the average flux:
\begin{equation}
    \begin{split}
        F_{i+1/2} &= \frac{1}{\Delta t} \int\limits_{t^{n}}^{t^{n+1}} f(u(x_{i+1/2},t)) \,\mathrm{d}t \\
                  &\approx \frac{1}{\Delta t} \int\limits_{t^{n}}^{t^{n+1}} g(u(x_{i+1/2},t)) \,\mathrm{d}t \\
                  &= \frac{1}{\Delta t} \int\limits_{t^{n}}^{t^{n+1}}
                  f(u_{i}^{n+1}) + a_i \left( u_i^+ - u_{i}^{n+1} - \sigma_i a_i \left( t - t^n \right) \right)
                  \,\mathrm{d}t \\
                  &= f(u_{i}^{n+1}) + a_i \left( u_i^+ - u_{i}^{n+1} \right) - \sigma_i a_i^2 \frac{\Delta t}{2}
    \end{split}
\end{equation}
Introducing the ratio \(\lambda = \frac{\Delta t}{\Delta x}\) and substituting the expression for \(u_i^+\), we arrive at the final second-order implicit nonlinear flux:
\begin{equation}
    \label{eq:second_order_nonlinear_flux}
    F_{i+1/2} \approx f(u_{i}^{n+1}) +
    a_i \sigma_i \Delta x \frac{1}{2} \left(
    1 + (2 f'(u_{i}^{n+1}) - a_i) \lambda
\right).
\end{equation}

By defining the local Courant number \(\nu_i = (2 f'(u_{i}^{n+1}) - a_i) \lambda\), the resulting flux is the direct nonlinear analogue of the linear implicit case:
\[
    F_{i+1/2} \approx f(u_{i}^{n+1}) +
    a_i \sigma_i \Delta x \frac{1}{2} \left(
    1 + \nu_i \right).
\]

\subsection*{Remark: Wave Speeds}%
\label{sub:wave_speeds}

It is worth noting a subtle but important distinction in the wave speeds used in the derivation above.
In the linear advection case, the speed \(a\) is a constant that simultaneously defines both the characteristic speed and the rate of flux change.
In the nonlinear case, these roles are decoupled to provide some flexibility.

Specifically, \(f'(u_{i}^{n+1})\) represents the \textit{characteristic speed} that determines how far the value \(u_{i}^{n+1}\) is shifted from \(u_{i}^{n}\).
On the other hand, \(a_i\) represents the \textit{linearized speed} used to approximate the flux function \(f(u)\) around the cell average.

While setting \(a_i = f'(u_{i}^{n+1})\) is the most direct approach, other choices are possible, as shown later.
% TODO: fig showing the different wave speeds
